\section{Contextualização}

\subsection{Cenário geral}
{Materiais compósitos}\\
Os materiais compósitos são amplamente utilizados em diversas indústrias devido às suas propriedades superiores, como alta resistência e leveza. Exemplos comuns incluem compósitos laminados, como fibra de vidro e fibra de carbono, que são empregados em setores como aeroespacial, automotivo e construção civil \cite{silva2019}.


{Compósitos micro-periódicos}\\
Os compósitos micro-periódicos são caracterizados por uma estrutura que varia rapidamente em uma escala muito pequena, embora possam parecer homogêneos a olho nu. Essa microestrutura permite obter propriedades físicas melhoradas em comparação com materiais homogêneos, tornando-os ideais para aplicações que exigem desempenho superior \cite{pereira2021}.

\subsection{Desafios na modelagem}
A modelagem de compósitos apresenta desafios significativos devido à complexidade de suas microestruturas e à interação entre as diferentes fases. Métodos tradicionais, como o método dos elementos finitos (MEF), podem ser utilizados para simular o comportamento mecânico dos compósitos, mas a precisão dos resultados é comprometida pela necessidade de um detalhamento muito fino da malha, o que aumenta o custo computacional \cite{oliveira2020}.

\subsection{MHA}

{Passos do MHA}\\
Os passos do MHA incluem a definição do problema, a escolha das escalas de homogeneização, a formulação das equações diferenciais e a aplicação de condições de contorno apropriadas. Em seguida, as soluções são obtidas para as escalas micro e macroscópicas, e as propriedades efetivas do material são calculadas a partir dessas soluções.

\subsection{MEF}
O Método dos Elementos Finitos (MEF) é uma técnica numérica amplamente utilizada para a resolução de problemas de engenharia e física. O MEF envolve a discretização do domínio do problema em elementos finitos, permitindo a aproximação das soluções das equações diferenciais que governam o comportamento do sistema. Essa abordagem é especialmente útil para problemas complexos com geometria irregular e condições de contorno complicadas. O MEF tem sido aplicado com sucesso em diversas áreas, incluindo mecânica estrutural, dinâmica de fluidos e transferência de calor.

{Passos do MEF}\\ 
Os passos do MEF incluem a discretização do domínio do problema em elementos finitos, a formulação das equações diferenciais em termos desses elementos, a aplicação de condições de contorno e a resolução do sistema de equações resultante. Em seguida, as soluções são obtidas para cada elemento e combinadas para obter a solução global do problema.

\subsection{IA para EDPs}

{PINNs}\\
As Redes Neurais de Física (PINNs) são uma abordagem inovadora que combina redes neurais com equações diferenciais parciais (EDPs) para resolver problemas de modelagem física. Essa técnica permite incorporar conhecimento prévio sobre o sistema em estudo, como leis da física e condições de contorno, diretamente na arquitetura da rede neural. Como resultado, as PINNs são capazes de aprender soluções para EDPs de forma mais eficiente e precisa do que métodos numéricos tradicionais. As PINNs têm sido aplicadas com sucesso em diversas áreas, incluindo mecânica dos fluidos, transferência de calor e elasticidade, demonstrando seu potencial para resolver problemas complexos de modelagem física.

